\documentclass{article} 
\input{../../lib.sty} 
\allowdisplaybreaks

\title{\texttt{fn function\_report\_noisy\_top\_k}} 
\author{Michael Shoemate} 
\begin{document}  
\maketitle 
 
\section{Hoare Triple} 
\subsection*{Precondition} 
\subsubsection*{Compiler-verified}
Types consistent with pseudocode.

\subsubsection*{Caller-verified}
\texttt{scale} is positive.
 
\subsection*{Pseudocode} 
\label{sec:python-pseudocode} 
\lstinputlisting[language=Python,firstline=2,escapechar=|]{./pseudocode/function_report_noisy_top_k.py} 
 
\subsection*{Postcondition} 

\begin{theorem}
    Returns a noninteractive function with no side-effects that, 
    when given a vector of non-null scores,
    returns the indices of the top k $z_i$,
    where each $z_i \sim RV(\mathrm{shift}=y_i, \mathrm{scale}=\texttt{scale})$,
    and each $y_i = -x_i$ if \texttt{optimize} is \texttt{min}, else $y_i = x_i$.

    The returned function will only return an error if entropy is exhausted.
    If an error is returned, the error is data-dependent.
\end{theorem}

\begin{proof}
    Assume the scores are non-null, as required by the returned function precondition.
    Therefore casts on line \ref{fn-cast} should never fail.
    However, if the input data is not in the input domain, and a score is null,
    then line \ref{fn-filter-nan} will filter out failed casts.
    This can be seen as a 1-stable transformation of the input data.

    The algorithm then proceeds to line \ref{fn-normalize-sign}.
    Assuming \texttt{optimize} is \texttt{max}, the line is a no-op, otherwise it negates each score,
    therefore each $y_i = -x_i$ if \texttt{optimize} is \texttt{min}, else $y_i = x_i$.

    The algorithm proceeds to line \ref{fn-init-sample}.
    The output measure has an associated noise distribution that is encoded into the Rust type system via
    \rustdoc{measurements/report\_noisy\_top\_k/trait}{SelectionMeasure}.
    \texttt{MO.random\_variable} on line \ref{fn-rv} creates a random variable \texttt{rv} 
    distributed according to \texttt{MO::RV} and parameterized by \texttt{shift} (the score), and \texttt{scale}.

    To sample from this random variable, line \ref{fn-partial-sample} constructs an instance of 
    \rustdoc{traits/samplers/psrn/struct}{PartialSample}, 
    which represents an infinitely precise sample from the random variable \texttt{rv}.
    We now have an iterator of pairs containing the index and noisy score of each candidate.

    The algorithm proceeds to line \ref{fn-max-sample},
    which defines a reducer based on \rustdoc{traits/samplers/psrn/struct}{PartialSample}\texttt{greater\_than}.
    Assume the scores are non-null, as required by the returned function precondition.
    Then \texttt{max\_sample} on line \ref{fn-max-sample} defines a total ordering on the scores.

    Since the score vector is finite, and \texttt{max\_sample} defines a total ordering,
    then the preconditions for \rustdoc{measurements/report_noisy_top_k/fn}{top} are met.
    Therefore on line \ref{fn-top} \texttt{top} returns the pairs with the top $k$ scores.
    Line \ref{fn-return-indices} then discards the scores, returning only the indices,
    which is the desired output.

    If entropy is exhausted, then the algorithm will return an error from \rustdoc{traits/samplers/psrn/struct}{PartialSample}\texttt{greater\_than}.
    Otherwise, there is one source of error in the function,
    when there are no non-null scores in the input vector.
\end{proof}

The algorithm avoids the pitfall of materializing an infinitely precise sample in memory by comparing finite arbitrary-precision bounds 
on the noisy scores until the lower bound of one noisy score is greater than the upper bound of all others.

\end{document} 
